% --------------------- NO TOCAR ----------------------
\renewcommand{\abstractname}{Resumen}
% -----------------------------------------------------

% HAZ UN RESUMEN DE TU TRABAJO. ES UN RESUMEN, NO ES UNA INTRODUCCIÓN
\begin{abstract}
    Este informe presenta un estudio comparativo de algoritmos de Aprendizaje por Refuerzo 
    estructurado en tres partes. 
    
    En la primera, se aborda el problema del bandido de $k$-brazos, 
    evaluando las familias $\epsilon$-Greedy, Softmax y UCB bajo distribuciones de recompensa Normal, 
    Bernoulli y Binomial. Los resultados muestran la superioridad de UCB1, que alcanza un 99\% de 
    selección del brazo óptimo con crecimiento logarítmico del \textit{regret}, frente a la 
    sensibilidad paramétrica de Softmax y la exploración residual de $\epsilon$-Greedy. 

    En la segunda y tercera parte, se extiende el estudio a entornos secuenciales modelados 
    como Procesos de Decisión de Markov. Sobre el entorno discreto \textbf{Taxi-v3}, se 
    implementan y comparan métodos tabulares: Montecarlo (on/off-policy), SARSA, Expected-SARSA, 
    Q-Learning y Double Q-Learning, confirmando la mayor eficiencia de los enfoques \textit{off-policy}. 
    Para el entorno continuo \textbf{LunarLander-v3}, se demuestra el fracaso de la discretización 
    mediante \textit{tile coding} y la necesidad de aproximadores neuronales: Deep SARSA logra un 
    89.7\% de aterrizajes exitosos, mientras que Deep Q-Learning (DQN), apoyado en 
    \textit{Experience Replay} y \textit{Target Network}, alcanza un 96.7\%, 
    consolidándose como el algoritmo más robusto y eficiente. 
    
    El lector podrá encontrar los experimentos realizados explicados y justificados en mayor detalle 
    en los notebooks que se han desarrollado. Toda la experimentación 
    es reproducible mediante estos notebooks, alojados en el repositorio de GitHub asociado a este proyecto \cite{repo2025carrillo_ibarrola_palomar}.
 \end{abstract}